\documentclass[10pt,twocolumn,letterpaper]{article}

\usepackage{amsmath,amssymb,booktabs,algorithm,algorithmic}
\usepackage{graphicx}

\title{Risk-Aware Windowed Conflict Replanning for Nonholonomic Multi-Robot Warehouse Navigation with Dynamic Pedestrians}
\author{Project Draft}
\date{}

\begin{document}
\maketitle

\begin{abstract}
Autonomous warehouse fleets must execute many transport tasks concurrently while remaining collision-free with shelves, other robots, and humans. Classical multi-agent path finding (MAPF) methods provide strong tools for robot-robot coordination, but their discrete assumptions and static obstacle models are not sufficient for human-shared warehouse spaces, where pedestrians move continuously and their future motion is uncertain. This paper presents a practical risk-aware multi-robot planning framework for nonholonomic warehouse robots operating around dynamic pedestrians. The proposed method combines kinodynamic lattice A*, Gaussian short-horizon pedestrian risk fields, time-expanded robot reservations, and bounded windowed conflict replanning. The system first computes dynamically feasible single-robot trajectories, then repairs local multi-robot conflicts in a finite future window while penalizing proximity to predicted pedestrian motion. The resulting trajectories are smoothed for presentation using continuous interpolation and corner rounding, while collision and near-miss metrics are evaluated on synchronized timed paths.

We validate the method on a 4-robot, 3-pedestrian warehouse scenario and additional larger environments with up to 6 robots, 5 pedestrians, and 12 static obstacles. Compared with independent A*, prioritized planning, risk-aware independent planning, risk-aware prioritized planning, and a CBS-style reservation proxy, the proposed method achieves zero robot/wall collisions and zero human near-threshold events in the main scenario. Ablations over human-risk weight, risk-field spread, and planning-window size show that dynamic pedestrian risk is the primary factor controlling clearance, while the bounded conflict window provides a scalable compromise between centralized optimal planning and purely sequential planning. The experiments highlight a safety-efficiency trade-off: human-aware trajectories are longer and more computationally expensive, but they eliminate high-risk human interactions that remain present in risk-unaware baselines.
\end{abstract}

\section{Introduction}
Multi-robot warehouse automation is increasingly important in modern logistics, where fleets of mobile robots transport shelves, parcels, and inventory between storage, packing, and loading areas. In such systems, planning is not only a shortest-path problem. A solution must be \emph{collision-free}: every robot trajectory must avoid static map geometry such as walls and shelves, avoid other robots in space-time, and maintain safe separation from dynamic agents such as workers or pedestrians. If robot $i$ has body radius $r_i$ and position $p_i(t) \in \mathbb{R}^2$, collision-free execution requires
\begin{equation}
    B(p_i(t), r_i) \cap \mathcal{O}_{\mathrm{static}} = \emptyset,
\end{equation}
\begin{equation}
    \|p_i(t) - p_k(t)\|_2 > r_i + r_k, \quad \forall i \neq k,
\end{equation}
and, for each pedestrian $h_j(t)$ with radius $r^h_j$,
\begin{equation}
    \|p_i(t) - h_j(t)\|_2 > r_i + r^h_j.
\end{equation}
In real deployments, however, the last condition is not enough. A robot that barely avoids a human can still be unsafe or socially unacceptable. We therefore evaluate not only hard collisions, but also near-threshold events under two margins, $\delta_{\mathrm{small}}=0.35$m and $\delta_{\mathrm{large}}=1.00$m, beyond the combined robot-human radii.

Robot-robot collision avoidance has been studied extensively through MAPF, prioritized planning, Conflict-Based Search (CBS), and bounded-suboptimal variants such as ECBS and EECBS. In a warehouse fleet, all robots are usually centrally controlled, so the system can reserve space-time resources, impose priorities, or replan conflicts. This makes robot-robot coordination challenging but structurally manageable. The more difficult problem is safe operation around people. Pedestrians are not controlled by the robot fleet, may deviate from nominal routes, and can enter aisles in ways that are only predictable over a short horizon. Consequently, a warehouse planner must handle two types of uncertainty at once: robot-robot interactions that can be scheduled, and human motion that must be anticipated but cannot be commanded.

This project focuses on that second challenge while preserving the first. We propose a risk-aware nonholonomic multi-robot planning framework that treats pedestrians as short-horizon dynamic risk sources and coordinates robots with a bounded conflict-window repair procedure. The method is deliberately practical: it does not attempt full optimal joint kinodynamic optimization over all robots and all times. Instead, it uses a kinodynamic A* lattice for each robot, evaluates human risk along motion primitives, reserves timed cells for robot-robot separation, and repairs only the portions of paths involved in near-future conflicts.

The contributions of this work are:
\begin{itemize}
    \item A nonholonomic multi-robot warehouse planner that integrates kinodynamic lattice search, dynamic pedestrian risk, time-expanded reservations, and bounded conflict-window replanning.
    \item A Gaussian pedestrian risk-field formulation using short-term predicted human positions at multiple temporal offsets.
    \item A practical coordination strategy that improves simultaneous robot motion without requiring full joint-state search.
    \item A complete experimental suite saved under \texttt{results/0615\_results/}, including algorithm comparisons, risk ablations, planning-window ablations, scenario diversity tests, heatmap path overlays, and paper-style plots.
\end{itemize}

\begin{figure}[t]
    \centering
    \caption{System overview: warehouse map, robots, pedestrians, static shelves, predicted human risk field, and planned nonholonomic trajectories.}
    \label{fig:system_overview}
\end{figure}

\section{Related Works}
\paragraph{Multi-Agent Path Finding.}
MAPF formulates the problem of moving multiple agents from starts to goals without inter-agent conflicts. A valid MAPF solution typically forbids vertex conflicts, edge-swap conflicts, and sometimes additional body-overlap conflicts. Optimal MAPF is computationally difficult, motivating both optimal and bounded-suboptimal solvers. CBS is a leading optimal approach that separates the problem into a high-level constraint tree and low-level single-agent path searches. When a collision is found, CBS branches by adding constraints that forbid one of the colliding agents from occupying the conflicting space-time resource. ECBS and EECBS accelerate this idea using focal search and inadmissible estimates, trading exact optimality for bounded-suboptimal runtime.

These methods are powerful for robot-robot coordination, but standard MAPF usually assumes grid-like holonomic agents, discrete time, and known static environments. Warehouse robots are nonholonomic, have finite turning radius, and may operate around continuously moving humans. Directly applying full CBS or ECBS with kinodynamic constraints can become expensive because each low-level search is itself a continuous or hybrid-state planning problem.

\paragraph{Windowed and Lifelong MAPF.}
Rolling-horizon methods, including Rolling-Horizon Collision Resolution (RHCR), decompose long or lifelong warehouse tasks into bounded time windows. Instead of resolving every future conflict at once, the solver resolves conflicts only in a local horizon and repeats the process as time advances. This is attractive for warehouses because tasks and obstacles change over time. Our method follows the same practical motivation, but uses a nonholonomic risk-aware lattice planner and a conflict-window repair mechanism rather than a full windowed MAPF solver.

\paragraph{Nonholonomic and Kinodynamic Multi-Robot Planning.}
Car-like MAPF and kinodynamic CBS extend MAPF ideas to robots with heading, steering, velocity, and body-shape constraints. CL-CBS uses spatiotemporal hybrid-state A* for car-like agents, while K-CBS adapts CBS to kinodynamic multi-robot motion planning in continuous domains. These methods are closer to real robot kinematics but are more computationally demanding than discrete MAPF. Our planner adopts the same principle that the low-level path should be dynamically feasible, but uses a lightweight lattice and bounded repair strategy to remain practical for repeated experiments and visualization.

\paragraph{Dynamic Obstacle and Human-Aware Navigation.}
Dynamic obstacle avoidance has also been studied through reactive velocity methods such as Dynamic Window Approach (DWA), velocity obstacles, and ORCA/RVO. These methods are fast and useful for local collision avoidance, but they are often formulated as reactive local controllers rather than global multi-robot task planners. Human-aware navigation additionally uses social-force models, risk maps, costmaps, and trajectory optimization methods such as Timed Elastic Bands (TEB) or nonlinear MPC. These approaches model proximity to people as a cost or constraint, but they usually focus on one robot or local navigation. In contrast, our system keeps the multi-robot warehouse planning structure while adding a human-risk field to the low-level kinodynamic search.

\paragraph{Remaining Gap.}
Existing methods tend to emphasize either multi-robot optimality or dynamic-obstacle reactivity. Full CBS/ECBS-style solvers provide strong robot-robot guarantees but become expensive under kinodynamic constraints and dynamic pedestrians. Local dynamic-obstacle methods react to humans but do not naturally solve centralized multi-robot scheduling. Our work targets the gap between these extremes: a centralized, risk-aware, nonholonomic planner that is less optimal than full CBS but easier to run, visualize, and extend for human-shared warehouses.

\section{Methodology}
\subsection{Problem Formulation}
Let $\mathcal{W} \subset \mathbb{R}^2$ be the warehouse workspace, $\mathcal{O}$ the set of static obstacles, $\mathcal{R}=\{1,\ldots,N\}$ the robots, and $\mathcal{H}=\{1,\ldots,M\}$ the pedestrians. Robot $i$ has state
\begin{equation}
    x_i(t) = [p^x_i(t), p^y_i(t), \theta_i(t)]^\top,
\end{equation}
start $s_i$, goal $g_i$, and body radius $r_i$. Pedestrian $j$ is represented by a time-indexed predicted position $\hat{h}_j(t)$ and radius $r^h_j$. The objective is to compute timed robot trajectories
\begin{equation}
    \Pi = \{\pi_i\}_{i=1}^{N}, \quad
    \pi_i = \{x_i^0, x_i^1, \ldots, x_i^{T_i}\},
\end{equation}
that reach the goals, avoid static obstacles, avoid robot-robot conflicts, and reduce exposure to predicted pedestrian motion.

We evaluate clearance between robot $i$ and pedestrian $j$ as
\begin{equation}
    c_{ij}(t) = \|p_i(t) - \hat{h}_j(t)\|_2 - (r_i + r^h_j).
\end{equation}
A hard robot-human collision occurs when $c_{ij}(t) \leq 0$. A near-threshold event occurs when $0 < c_{ij}(t) \leq \delta$, where we report both $\delta=0.35$m and $\delta=1.00$m.

\subsection{Algorithm Overview}
The proposed planner, which we call \emph{Risk-Aware Windowed Conflict Replanning}, consists of five stages:
\begin{enumerate}
    \item Predict pedestrian motion over a short time horizon and construct a dynamic risk cost.
    \item Plan nonholonomic single-robot paths with kinodynamic lattice A*.
    \item Build time-expanded reservations from previously planned robot paths.
    \item Scan a bounded future window for robot-robot conflicts.
    \item Replan only the lower-priority robot suffix around the conflict region.
\end{enumerate}

\begin{algorithm}[t]
\caption{Risk-Aware Windowed Conflict Replanning}
\label{alg:rawcr}
\begin{algorithmic}[1]
\STATE \textbf{Input:} map $\mathcal{W}$, robots $\mathcal{R}$, pedestrians $\mathcal{H}$, window length $H$, repair iterations $K$
\STATE Predict pedestrian positions $\hat{h}_j(t+\tau)$ for short offsets $\tau \in \mathcal{T}$
\STATE Build a risk-aware kinodynamic lattice planner $A^*_{\mathrm{risk}}$
\STATE Compute initial paths $\Pi$ using prioritized planning and time-expanded reservations
\FOR{$k=1$ to $K$}
    \STATE Scan $\Pi$ for the first robot-robot conflict within $[0,H\Delta t]$
    \IF{no conflict exists}
        \STATE \textbf{break}
    \ENDIF
    \STATE Select the lower-priority robot $i$ involved in the conflict
    \STATE Set replan time $t_r = \max(0, t_{\mathrm{conflict}} - L\Delta t)$
    \STATE Reserve all other robot paths relative to $t_r$
    \STATE Replan the suffix of $\pi_i$ from $x_i(t_r)$ to $g_i$ using $A^*_{\mathrm{risk}}$
    \STATE Replace the old suffix if replanning succeeds
\ENDFOR
\STATE \textbf{Return:} timed robot trajectories $\Pi$
\end{algorithmic}
\end{algorithm}

\subsection{Nonholonomic Kinodynamic Lattice}
Each robot is planned in a discrete hybrid state space
\begin{equation}
    q = (u, v, b, n),
\end{equation}
where $(u,v)$ are grid indices, $b$ is a heading bin, and $n$ is a time index. The continuous heading corresponding to bin $b$ is
\begin{equation}
    \theta(b) = \frac{2\pi b}{B},
\end{equation}
where $B$ is the number of heading bins. The planner uses a small set of motion primitives:
\begin{equation}
    \mathcal{U} = \{\mathrm{forward}, \mathrm{left}, \mathrm{right}, \mathrm{wait}\}.
\end{equation}
For a primitive with distance $d$ and heading increment $\Delta b$, the next state is
\begin{equation}
    b' = (b+\Delta b) \bmod B,
\end{equation}
\begin{equation}
    p' = p + d[\cos\theta(b'), \sin\theta(b')]^\top,
\end{equation}
\begin{equation}
    t' = t + \Delta t.
\end{equation}
The wait primitive has $d=0$ and increments time only. This lattice is not a full continuous bicycle model optimization, but it encodes finite turning behavior, heading continuity, and time-indexed occupancy. Transitions are sampled at sub-cell resolution to reject static and dynamic collisions along the segment, not only at the endpoints.

\subsection{Dynamic Pedestrian Risk Field}
Pedestrians are represented as moving circular dynamic obstacles. For planning, the system predicts pedestrian positions at a small set of temporal offsets
\begin{equation}
    \mathcal{T} = \{-0.6, 0.0, 0.6, 1.2\} \ \mathrm{s},
\end{equation}
which makes the robot sensitive not only to exact-time collision but also to near-future pedestrian occupancy. For robot state $x_i(t)$, the clearance to pedestrian $j$ at offset $\tau$ is
\begin{equation}
    c_{ij}^{\tau}(t) =
    \|p_i(t)-\hat{h}_j(t+\tau)\|_2 - (r_i+r^h_j).
\end{equation}
If $c_{ij}^{\tau}(t) \leq 0$, the transition is assigned infinite cost because it intersects a dynamic obstacle. If the clearance is positive but below a safety radius $d_{\mathrm{safe}}$, the planner adds a Gaussian risk cost
\begin{equation}
    R_i(t) =
    \sum_{j=1}^{M}\sum_{\tau\in\mathcal{T}}
    \lambda_h \gamma(\tau)
    \exp\left(
    -\frac{1}{2}\left(\frac{c_{ij}^{\tau}(t)}{\sigma}\right)^2
    \right),
\end{equation}
where $\lambda_h$ is the human-risk weight, $\sigma$ is the risk-field spread, and $\gamma(\tau)$ is a temporal decay factor. In the implementation, $\gamma(0)=1$ and nonzero offsets receive a smaller decay weight. This risk is averaged over sampled points along a motion primitive, producing smooth cost pressure away from predicted pedestrian paths.

\begin{figure}[t]
    \centering
    \caption{Human-risk heatmap and risk-weight path overlay. The same robot/environment is shown with multiple risk weights, where lighter paths correspond to smaller risk weights and darker paths correspond to larger risk weights.}
    \label{fig:risk_overlay}
\end{figure}

\subsection{Single-Robot Search Objective}
The low-level planner minimizes a weighted path cost
\begin{equation}
    J_i(\pi_i) =
    \sum_{n=0}^{T_i-1}
    \left[
    w_d d_n + w_{\theta}|\Delta \theta_n|
    + w_{\mathrm{wait}}\mathbb{I}_{\mathrm{wait}}
    + R_i(t_n)
    \right]
    + w_g h(x_i^{T_i}, g_i),
\end{equation}
where $d_n$ is primitive distance, $\Delta\theta_n$ is heading change, $\mathbb{I}_{\mathrm{wait}}$ indicates a wait action, $R_i(t_n)$ is the pedestrian risk term, and $h$ is a Euclidean distance-to-goal heuristic. This objective encourages short paths, smooth heading changes, limited waiting, and larger pedestrian clearance.

\subsection{Time-Expanded Reservations}
For robot-robot coordination, the planner builds a reservation table from previously planned robot paths. Let $q_i^n$ be the grid cell occupied by robot $i$ at time index $n$. A vertex reservation prevents later robots from occupying the same inflated cell:
\begin{equation}
    (q_i^n, n) \in \mathcal{V}_{\mathrm{res}}.
\end{equation}
An edge-swap reservation prevents two robots from exchanging cells over the same time step:
\begin{equation}
    (q_i^n, q_i^{n+1}, n) \in \mathcal{E}_{\mathrm{res}}.
\end{equation}
The reservation table can be inflated by a cell-padding parameter to create a spatial buffer around planned robots. This is similar in spirit to prioritized planning, but our proposed method subsequently scans and repairs conflicts in a bounded window rather than accepting the first sequential solution.

\subsection{Bounded Conflict-Window Replanning}
After initial prioritized planning, the proposed method scans a finite future window of length
\begin{equation}
    T_H = H\Delta t.
\end{equation}
In the main experiments, $H=32$ and $\Delta t=0.2$s, so the conflict window covers $6.4$s. A conflict is detected if two robots become closer than their combined radii plus a clearance margin:
\begin{equation}
    \|p_i(t)-p_k(t)\|_2
    < r_i + r_k + \epsilon_{\mathrm{rr}}.
\end{equation}
When a conflict is found, the method chooses the lower-priority robot, keeps its prefix up to $t_r$, and replans its suffix while treating all other robots as time-expanded reservations. This approximates windowed MAPF and CBS-style repair, but avoids constructing a full constraint tree. It is therefore heuristic rather than optimal, but it is easier to run on nonholonomic paths with dynamic pedestrian costs.

\subsection{Receding-Horizon Interpretation}
The implemented proposed planner is best understood as \emph{receding-horizon-inspired} rather than a full model predictive controller. It does not jointly optimize all controls at every execution step. Instead, it applies a bounded lookahead window for conflict detection and local suffix repair. This is appropriate for dynamic warehouses because conflicts far in the future may be invalidated by new human motion or replanning. The project also contains a rolling-window planner variant, but the reported experiments use the faster bounded conflict-window replanner.

\subsection{Trajectory Smoothing and Visualization}
The kinodynamic planner returns lattice paths. For visualization, paths are rounded using open Chaikin subdivision. Given adjacent points $P_k$ and $P_{k+1}$, the smoother inserts
\begin{equation}
    Q_k = (1-\alpha)P_k + \alpha P_{k+1},
\end{equation}
\begin{equation}
    R_k = \alpha P_k + (1-\alpha)P_{k+1},
\end{equation}
with $\alpha \in [0,0.5)$. Repeated subdivision removes sharp visual corners while preserving the route topology. Dense time interpolation then produces continuous robot poses for paper-style videos and figures. Importantly, quantitative metrics are computed on timed planner paths, while smoothing is used for presentation-quality rendering.

\subsection{Metrics}
We evaluate:
\begin{itemize}
    \item \textbf{Total distance}: $\sum_i \sum_n \|p_i^{n+1}-p_i^n\|_2$.
    \item \textbf{Makespan}: $\max_i T_i$.
    \item \textbf{Computation time}: wall-clock planning runtime.
    \item \textbf{Wall collision count}: sampled robot footprint collisions with static occupancy.
    \item \textbf{Robot collision count}: sampled robot-robot footprint overlaps.
    \item \textbf{Human near count}: number of samples with robot-human clearance under $0.35$m and $1.00$m.
    \item \textbf{Minimum human clearance}: $\min_{i,j,t} c_{ij}(t)$.
\end{itemize}

\section{Experiments}
\subsection{Implementation and Output Directory}
All reported experiments are generated by
\begin{equation}
    \texttt{scripts/run\_0615\_experiments.py}.
\end{equation}
The outputs are saved under
\begin{equation}
    \texttt{results/0615\_results/}.
\end{equation}
The main scenario contains 4 robots, 3 pedestrians, and 10 static shelves in an $18$m by $12$m warehouse. Additional diversity scenarios contain 5 robots/4 pedestrians/8 obstacles and 6 robots/5 pedestrians/12 obstacles. Pedestrian waypoint intents are converted to obstacle-aware pedestrian rollouts before being used as dynamic obstacles, preventing pedestrians from passing through shelves in generated visualizations.

\subsection{Experiment 1: Algorithm Comparison}
We compare six planners:
\begin{enumerate}
    \item Independent A*: no robot coordination and no human-risk cost.
    \item Prioritized Planning: robot reservations but no human-risk cost.
    \item Risk-aware Independent A*: human-risk cost but no robot coordination.
    \item Risk-aware Prioritized: human-risk cost plus sequential reservations.
    \item CBS-style Reservations: inflated reservation proxy, not full CBS.
    \item Proposed Windowed Risk-Aware: risk-aware lattice A* plus bounded conflict-window repair.
\end{enumerate}

\begin{table}[t]
    \centering
    \caption{Algorithm comparison on the 4-robot / 3-pedestrian warehouse scenario. Insert \texttt{results/0615\_results/algorithm\_comparison.csv} or the rendered table image here.}
    \label{tab:algorithm_comparison}
\end{table}

\begin{figure}[t]
    \centering
    \caption{Multi-panel algorithm comparison: total distance, computation time, wall/robot collision count, human near count at $0.35$m and $1.00$m, and minimum human clearance.}
    \label{fig:algorithm_comparison}
\end{figure}

The risk-unaware baselines produce the shortest trajectories but violate safety. Independent A* and Prioritized Planning both travel $79.39$m with makespan $7.6$s, but they produce 6 robot-robot collisions, 9 small-threshold human near events, 22 large-threshold human near events, and a negative minimum human clearance of $-0.322$m. This indicates an actual human collision or overlap under the dynamic-obstacle model.

The CBS-style reservation proxy removes robot collisions but does not account for human risk. It reports zero robot collisions, but still produces 6 small-threshold and 20 large-threshold human near events, with the same negative minimum human clearance. This result is important: solving robot-robot coordination alone does not solve human-aware navigation.

Risk-aware Independent A* eliminates human near events and increases minimum human clearance to $1.260$m, but because it does not coordinate robots, it produces 10 robot collisions. Risk-aware Prioritized Planning removes both robot collisions and human near events, but requires $60.28$s of computation. The proposed Windowed Risk-Aware method also achieves zero wall/robot collisions and zero human near events, with the same $1.260$m minimum human clearance, while reducing computation time to $46.03$s. Relative to risk-aware prioritized planning, the proposed method is approximately
\begin{equation}
    \frac{60.28 - 46.03}{60.28} \times 100 \approx 23.6\%
\end{equation}
faster on the main scenario.

The safety improvement comes with a path-efficiency cost. The proposed method travels $95.15$m, about $19.8\%$ longer than the risk-unaware independent baseline. This is the central trade-off in the project: safe human-aware trajectories are longer and more expensive, but they avoid unsafe pedestrian interactions.

\subsection{Experiment 2: Human Risk Ablation}
We vary the human-risk weight $\lambda_h \in \{0,2,4,8,12\}$ and risk spread $\sigma \in \{0.45,0.75,1.05,1.35\}$ in a fixed human-crossing scenario. The purpose is to isolate how the risk field changes the path of a single robot.

\begin{table}[t]
    \centering
    \caption{Risk-field ablation over human-risk weight and Gaussian risk spread. Insert \texttt{results/0615\_results/risk\_ablation\_metrics.csv}.}
    \label{tab:risk_ablation}
\end{table}

\begin{figure}[t]
    \centering
    \caption{Risk ablation curves showing how distance, minimum human clearance, and human near count change with risk weight and risk sigma.}
    \label{fig:risk_ablation}
\end{figure}

With $\lambda_h=0$, the path length is $7.13$m and the minimum human clearance is $0.765$m, but there are 3 large-threshold human near events. Increasing the risk weight to $\lambda_h=2$ removes all near events and raises minimum clearance to $1.481$m. At $\lambda_h=8$, the minimum clearance reaches $1.848$m, while the path length increases to $8.22$m. This confirms that the risk term has a visible geometric effect: it pushes the robot away from the predicted pedestrian trajectory rather than merely avoiding hard collision.

The sigma ablation shows that a narrow field, $\sigma=0.45$, already removes near events with $1.481$m minimum clearance. Increasing to $\sigma=0.75$ produces a wider avoidance arc with $1.848$m clearance and a longer path. Larger $\sigma$ values saturate in this scenario because the best available avoidance route is already selected.

\subsection{Experiment 3: Heatmap Path Overlay}
The heatmap visualization overlays multiple paths on the same environment. The risk field is shown as a smooth heatmap around predicted pedestrian positions, and paths become darker as the risk weight increases. This figure is designed to show that the planner is not only reporting better metrics but actually changing path geometry in response to human risk.

\begin{figure}[t]
    \centering
    \caption{Human-risk heatmap with overlaid paths for risk weights $0,2,4,8,12$. Use \texttt{results/0615\_results/risk\_heatmap\_path\_overlay.png}.}
    \label{fig:heatmap_overlay}
\end{figure}

\subsection{Experiment 4: Planning-Window Ablation}
We vary the bounded conflict-window length $H \in \{16,24,32,40\}$ steps. With $\Delta t=0.2$s, these correspond to $3.2$s, $4.8$s, $6.4$s, and $8.0$s.

\begin{table}[t]
    \centering
    \caption{Planning-window ablation. Insert \texttt{results/0615\_results/planning\_window\_ablation.csv}.}
    \label{tab:window_ablation}
\end{table}

\begin{figure}[t]
    \centering
    \caption{Planning-window ablation curves for distance, makespan, and wall/robot collision count.}
    \label{fig:window_ablation}
\end{figure}

In the main scenario, all tested windows achieve the same distance, makespan, and zero collision counts. Computation time remains around $46$s. This suggests that the first relevant conflicts in the scenario are already visible within a $3.2$s window, and longer windows do not change the repair decisions. This is a useful but limited result: it shows robustness for the current layout, but not that window size is irrelevant in general. More adversarial bottleneck scenarios would likely expose stronger differences.

\subsection{Experiment 5: Scenario Diversity}
We evaluate the proposed planner on three environments:
\begin{itemize}
    \item 4 robots, 3 pedestrians, 10 obstacles, $18$m by $12$m.
    \item 5 robots, 4 pedestrians, 8 obstacles, $14$m by $8.6$m.
    \item 6 robots, 5 pedestrians, 12 obstacles, $18$m by $12$m.
\end{itemize}

\begin{table}[t]
    \centering
    \caption{Scenario diversity results. Insert \texttt{results/0615\_results/scenario\_diversity\_metrics.csv}.}
    \label{tab:scenario_diversity}
\end{table}

\begin{figure}[t]
    \centering
    \caption{Scenario overview for the 4R/3H, 5R/4H, and 6R/5H warehouse maps.}
    \label{fig:scenario_overview}
\end{figure}

\begin{figure}[t]
    \centering
    \caption{Scenario diversity metrics: distance, computation time, and human near count across larger maps.}
    \label{fig:scenario_diversity}
\end{figure}

The proposed method succeeds on all three scenarios and reports zero wall/robot collision counts. In the main 4R/3H warehouse, it also reports zero human near events with $1.260$m minimum human clearance. In the denser 5R/4H and 6R/5H settings, the planner still avoids hard collisions but allows a small number of near-threshold events: 2 small and 3 large events in the medium map, and 2 small and 3 large events in the large map. Minimum clearance drops to $0.138$m and $0.162$m respectively. This indicates that as human and robot density increases, the planner preserves collision-free execution but may still pass close to humans when the map offers limited alternatives.

Computation time is not monotonic with map size. The 4R/3H case takes $46.03$s, the 5R/4H case $8.31$s, and the 6R/5H case $23.37$s. This occurs because runtime depends not only on the number of agents but also on obstacle placement, goal geometry, reservation conflicts, and the shape of the dynamic risk field. Dense maps with structured lanes can be easier than open maps where many alternative risky branches must be evaluated.

\section{Conclusions}
This paper presented a practical risk-aware planning framework for nonholonomic multi-robot warehouse navigation with dynamic pedestrians. The key result is that robot-robot coordination alone is insufficient in human-shared environments. In the main experiment, a CBS-style reservation proxy removed robot collisions but still produced 20 large-threshold human near events and negative minimum human clearance. Conversely, risk-aware independent planning removed human near events but produced 10 robot collisions. Safe warehouse execution required both components: human-risk-aware low-level planning and multi-robot coordination.

The proposed Windowed Risk-Aware planner achieved zero wall/robot collisions and zero human near events on the main 4R/3H scenario. Compared with risk-aware prioritized planning, it obtained the same safety outcome with about $23.6\%$ lower computation time. The planner is therefore a useful middle ground between simple prioritized planning and full optimal CBS/ECBS. It is not globally optimal, but it is interpretable, modular, and compatible with nonholonomic motion and dynamic pedestrian risk.

The risk ablations show a clear safety-efficiency trade-off. Increasing the human-risk weight raises minimum pedestrian clearance from $0.765$m to $1.848$m, but also increases path length from $7.13$m to $8.22$m in the ablation scenario. This is expected: safer paths route around predicted pedestrian motion and therefore sacrifice distance optimality. Similarly, wider Gaussian risk fields increase the planner's willingness to bend away from humans, but eventually saturate once the selected route is already separated.

The planning-window ablation shows stable performance for $H=16$ to $H=40$ in the main scenario. This suggests that a moderate $6.4$s window is sufficient for the current warehouse layout. However, this should not be overinterpreted. More difficult bottleneck environments may require longer windows or stronger high-level repair. The diversity experiments further show that the proposed method scales to larger scenarios with zero wall/robot collisions, but dense human traffic can still reduce clearance. Thus, future work should consider adaptive risk weights, learned pedestrian prediction, and stronger local replanning around human bottlenecks.

The current method has several limitations. First, it is a heuristic approximation, not full optimal CBS/ECBS and not full joint kinodynamic optimization. Second, pedestrian prediction is short-horizon and deterministic; it does not model multimodal human intent. Third, smoothing is used for visualization rather than as a certified dynamically feasible trajectory optimization layer. Fourth, ORCA/RVO and TEB/NMPC baselines were not included as direct comparisons because they operate mainly as local velocity or trajectory optimizers, while this project uses a global nonholonomic lattice planner. A future version should add these baselines with a fair nonholonomic projection layer.

Despite these limitations, the project demonstrates a useful research direction: centralized warehouse fleets can use strong robot-robot coordination, but safe deployment in human-shared environments requires dynamic risk awareness. The proposed planner provides a compact and extensible bridge between MAPF-style coordination and human-aware navigation.

\section{Reference}
\begin{thebibliography}{99}

\bibitem{hart1968astar}
P. E. Hart, N. J. Nilsson, and B. Raphael.
A formal basis for the heuristic determination of minimum cost paths.
\emph{IEEE Transactions on Systems Science and Cybernetics}, 1968.

\bibitem{sharon2015cbs}
G. Sharon, R. Stern, A. Felner, and N. R. Sturtevant.
Conflict-based search for optimal multi-agent pathfinding.
\emph{Artificial Intelligence}, 2015.

\bibitem{barer2014ecbs}
M. Barer, G. Sharon, R. Stern, and A. Felner.
Suboptimal variants of the conflict-based search algorithm for the multi-agent pathfinding problem.
\emph{Proceedings of the International Symposium on Combinatorial Search}, 2014.

\bibitem{li2021eecbs}
J. Li, W. Ruml, and S. Koenig.
EECBS: A bounded-suboptimal search for multi-agent path finding.
\emph{AAAI Conference on Artificial Intelligence}, 2021.

\bibitem{li2020rhcr}
J. Li, A. Tinka, S. Kiesel, J. W. Durham, T. K. S. Kumar, and S. Koenig.
Lifelong multi-agent path finding in large-scale warehouses.
\emph{arXiv:2005.07371}, 2020.

\bibitem{wen2020clmapf}
L. Wen, Z. Zhang, Z. Chen, X. Zhao, and Y. Liu.
CL-MAPF: Multi-agent path finding for car-like robots with kinematic and spatiotemporal constraints.
\emph{arXiv:2011.00441}, 2020.

\bibitem{kottinger2022kcbs}
J. Kottinger, S. Almagor, and M. Lahijanian.
Conflict-based search for multi-robot motion planning with kinodynamic constraints.
\emph{arXiv:2207.00576}, 2022.

\bibitem{moldagalieva2023dbcbs}
A. Moldagalieva, J. Ortiz-Haro, M. Toussaint, and W. Hoenig.
db-CBS: Discontinuity-bounded conflict-based search for multi-robot kinodynamic motion planning.
\emph{arXiv:2309.16445}, 2023.

\bibitem{van2011orca}
J. van den Berg, S. J. Guy, M. Lin, and D. Manocha.
Reciprocal n-body collision avoidance.
\emph{Robotics Research}, 2011.

\bibitem{fox1997dwa}
D. Fox, W. Burgard, and S. Thrun.
The dynamic window approach to collision avoidance.
\emph{IEEE Robotics and Automation Magazine}, 1997.

\bibitem{helbing1995social}
D. Helbing and P. Molnar.
Social force model for pedestrian dynamics.
\emph{Physical Review E}, 1995.

\bibitem{cai2020survey}
K. Cai, C. Wang, J. Cheng, C. W. de Silva, and M. Q.-H. Meng.
Mobile robot path planning in dynamic environments: A survey.
\emph{arXiv:2006.14195}, 2020.

\bibitem{dolgov2008hybrid}
D. Dolgov, S. Thrun, M. Montemerlo, and J. Diebel.
Practical search techniques in path planning for autonomous driving.
\emph{AAAI Workshop on Search in Artificial Intelligence and Robotics}, 2008.

\bibitem{rosmann2017teb}
C. Rosmann, F. Hoffmann, and T. Bertram.
Integrated online trajectory planning and optimization in distinctive topologies.
\emph{Robotics and Autonomous Systems}, 2017.

\bibitem{kulathunga2024rteb}
G. Kulathunga, A. Yilmaz, Z. Huang, I. Hroob, H. Arunachalam, L. Guevara, A. Klimchik, G. Cielniak, and M. Hanheide.
Resilient timed elastic band planner for collision-free navigation in unknown environments.
\emph{arXiv:2412.03174}, 2024.

\end{thebibliography}

\end{document}
